\begin{textsample}{NEG Dim 4 – 2012\_c5 – Score -1.00 – 2012\_c5/enade\_2012\_tec\_gest\_fin\_8526ca4e82.txt}  \label{ex:f4_neg_105}
CAPÍTULO 5 DISTRIBUIÇÃO DOS CONCEITOS 5.1 PANORAMA NACIONAL DA DISTRIBUIÇÃO DOS CONCEITOS A Tabela 5.1 apresenta a quantidade e distribuição de cursos de Tecnologia em Gestão Financeira participantes do ENADE/2012, por faixa de conceito e Grande Região.
A diferença entre os cursos tabulados neste capítulo e no capítulo 2 são os cursos sem conceito, em princípio, aqueles sem alunos concluintes que participassem da prova ou que tivessem tão somente um aluno concluinte11.
Observando -se os dados da Tabela 5.1, nota -se que, dos 166 cursos participantes, 71 ( 42,8 \% ) classificaram -se com conceito 3, o valor modal.
Este foi também o conceito modal nas regiões Nordeste ( 38,9 \% ), Sudeste ( 50,5 \% ) e Centro-Oeste ( 33,3 \% ).
Na regiões Norte e Sul os conceitos modais foram, respectivamente, 2 ( 40,0 \% ) e 4 ( 40,5 \% ).
O conceito 2 foi o segundo mais frequente em nível nacional ( 22,3 \%, correspondendo a 37 cursos ) e o conceito 4, o terceiro ( 21,7 \%, correspondendo a 36 cursos ).
Houve, ainda, 12 cursos ( 7,2 \% ) que receberam conceito 5 e apenas quatro cursos que receberam conceito 1 ( 2,4 \% ).
Seis dos 166 cursos de Tecnologia em Gestão Financeira ficaram sem conceito ( SC ).
A região Norte participou com cinco cursos ou 3,0 \% do total nacional.
Nesta região, o conceito 2 foi atribuído a dois cursos dos cinco participantes, o que equivale a 40,0 \% do total regional, como já comentado.
Dois cursos ficaram sem conceito ( SC ) e o outro curso foi avaliado com conceito 3.
A região Nordeste participou com 18 cursos ou 10,8 \% do total nacional.
Como já conceito modal para a região.
Os conceitos 2 e 4 foram atribuídos a, respectivamente, um e cinco cursos ( 5,6 \% e 27,8 \% ).
Os conceitos 1 e 5 foram atribuídos a, respectivamente, um e três cursos ( 5,6 \% e 16,7 \% ).
Nessa região, apenas um dos cursos ficou sem conceito ( SC ).
Dos 97 cursos participantes da região Sudeste, como já comentado, 49 ( 50,5 \% ) obtiveram conceito 3, o conceito modal.
O conceito 1 foi atribuído a três cursos ( 3,1 \% ), o conceito 2 a 25 ( 25,8 \% ), o conceito 4 a 14 cursos ( 14,4 \% ) e o conceito 5, a cinco cursos ( 5,2 \% ).
Apenas um curso ficou sem conceito ( SC ).
A região Sul contou com 37 cursos ou 22,3 \% do total nacional.
A predominância do conceito 4, como já comentado, foi de 40,5 \%, correspondentes a 15 dos 37 cursos participantes na região Sul.
O conceito 2 foi atribuído a sete cursos ( 18,9 \% ), o conceito 3 a 11 cursos ( 29,7 \% ) e o conceito 5, a três deles ( 8,1 \% ).
Apenas um dos cursos da região Sul ficou sem conceito ( SC ), enquanto nenhum recebeu conceito 1.
Como já comentado, dos nove cursos participantes na região Centro-Oeste, três ( 33,3 \% ) receberam conceito 3, o conceito modal.
Os demais cursos foram avaliados com conceito 2 ( dois cursos, 22,2 \% ), conceito 4 ( também dois cursos ) e conceito 5 ( um curso, 11,1 \% ), enquanto apenas um curso ficou sem conceito ( SC ) e nenhum foi \textbf{alocado} no conceito 1. 5.2 CONCEITOS POR CATEGORIA ADMINISTRATIVA E POR GRANDE REGIÃO A Tabela 5.2 apresenta a distribuição dos cursos participantes do ENADE/2012 de Tecnologia em Gestão Financeira, por Categoria Administrativa, de acordo com os conceitos por eles alcançados, segundo as Grandes Regiões brasileiras.
Dos 166 cursos participantes, cinco ( 3,0 \% ) eram ministrados em instituições públicas e 161 ( 97,0 \% ), em privadas.
De acordo com as informações da Tabela 5.2, em termos nacionais, entre instituições públicas e privadas : das 12 IES que receberam conceito 5, uma era pública e 11 eram privadas.
Dos cinco cursos participantes de IES públicas, dois foram atribuídos ao conceito 3, outros dois ao conceito 4 e um ao conceito 5.
Nesta categoria, nenhum dos cursos ficou sem conceito ( SC ).
Na rede privada, o conceito modal foi 3, com 69 cursos dos 161 da categoria.
Nesta categoria administrativa, apenas seis dos cursos ficaram sem conceito ( SC ).
Entre os demais cursos participantes, quatro receberam conceito 1 e 37, conceito 2.
O conceito 4 foi atribuído a 34 cursos e o conceito 5, como já comentado, a 11 cursos.
Na análise por região, observa -se que, na região Norte, as instituições privadas concentraram todos os cinco cursos participantes, dos quais dois ficaram sem conceito ( SC ), outros dois receberam conceito 2 e um curso, recebeu conceito 3.
As instituições públicas, como já comentado, não participaram com cursos nesta região.
Na região Nordeste, a rede privada concentrou todos os 18 cursos participantes.
Nesta Categoria Administrativa na região, predominaram os cursos com conceito 3, sete cursos correspondendo a 38,9 \% dos cursos oferecidos por IES privadas no Nordeste.
Apenas um curso, nesta combinação de região e categoria, ficou sem conceito ( SC ).
Os demais foram avaliados com conceitos 1 ( um curso correspondendo a 5,6 \% ), 2 ( também um curso ), 4 ( cinco cursos correspondendo a 27,8 \% ) e 5 ( três cursos correspondendo a 16,7 \% ).
As instituições públicas dessa região, como já comentado, não participaram com cursos.
Na região Sudeste, a proporção de cursos da rede privada foi de 94,8 \%, correspondendo a 92 dos 97 cursos participantes.
Nesta categoria, na região Sudeste, o conceito modal foi 3 ( 47 cursos ), enquanto apenas um curso ficou sem conceito ( SC ).
Os demais foram avaliados com conceito 1 ( três cursos ), conceito 2 ( 25 cursos ), conceito 4 ( 12 cursos ) e conceito 5 ( quatro cursos ).
Entre os cinco cursos em instituições públicas na região Sudeste, dois receberam conceito 3, outros dois, conceito 4 e um conceito 5.
Na região Sudeste, nesta categoria, nenhum curso ficou sem conceito ( SC ).
As instituições privadas concentraram todos os 37 cursos participantes da região Sul.
Desses, 15 ficaram com conceito 4, o conceito modal.
Nesta combinação de Categoria Administrativa e Grande Região, apenas um curso ficou sem receber conceito ( SC ).
Os demais foram avaliados com conceito 2 ( sete cursos ), conceito 3 ( 11 cursos ) e conceito 5 ( três cursos ).
As instituições públicas na região Sul, como já comentado, não participaram com cursos.
Na região Centro-Oeste, todos os nove cursos participantes eram de instituições privadas.
Destes, três concentraram -se no conceito 3, conceito modal.
Apenas um curso ficou sem conceito ( SC ) nesta categoria, enquanto nenhum foi \textbf{alocado} ao conceito 1.
Os demais receberam conceito 2 ( dois cursos ), conceito 4 ( outros dois cursos ) e conceito 5 ( um curso ).
As instituições públicas, nesta região, não participaram com nenhum curso. 5.3 CONCEITOS POR ORGANIZAÇÃO ACADÊMICA E POR GRANDE REGIÃO Na Tabela 5.3 encontra -se a distribuição dos conceitos atribuídos aos cursos participantes do ENADE/2012 na Área de Tecnologia em Gestão Financeira, por Organização Acadêmica, segundo as Grandes Regiões brasileiras.
Dos 166 cursos participantes, 47 eram oferecidos em Universidades, 27 em Centros Universitários e os demais 92 em Faculdades.
Esta distribuição corresponde a, respectivamente, 28,3 \%, 16,3 \% e 55,4 \% dos cursos.
De acordo com os dados apresentados, de todos os 12 cursos avaliados com conceito 5, cinco eram vinculados a Universidades.
Somente um curso em Universidades ficou sem conceito ( SC ), enquanto nenhum recebeu conceito 1.
Esse tipo de Organização Acadêmica teve o conceito 3 como modal, com 21 cursos ( 44,7 \% ).
Os demais cursos de Universidades avaliados receberam os conceitos 2 ( nove cursos ), 4 ( 11 cursos ) e conceito 5 ( cinco cursos, como já mencionado ).
Entre os cursos em Centros Universitários, o conceito modal também foi 3, com 12 cursos ( 44,4 \% ).
Neste tipo de organização acadêmica nenhum curso ficou sem conceito ( SC ) ou recebeu o conceito 1.
Os outros cursos neste tipo de Organização Acadêmica receberam os conceitos 2 ( oito cursos ), 4 ( seis cursos ) e 5 ( um curso ).
Nas Faculdades, cinco dos 92 cursos ficaram sem conceito ( SC ) e 38 ( 41,3 \% ) receberam o mesmo conceito modal dos demais tipos de Organização Acadêmica, conceito 3.
Dos demais cursos neste tipo de Organização Acadêmica, quatro receberam conceito 1, 20 conceito 2, 19 conceito 4 e seis cursos, o conceito 5.
Considerando -se separadamente as regiões brasileiras, verifica -se que, na região Norte, as Universidades concentraram somente um dos cinco cursos participantes, o qual foi avaliado com conceito 3.
Os Centros Universitários da região Norte também foram representados por somente um curso, o qual foi avaliado com conceito 2.
As Faculdades participaram com três cursos na região Norte.
Dois cursos ficaram sem conceito ( SC ) e o outro recebeu o conceito 2.
Na região Nordeste, as Universidades participaram com dois dos 18 cursos na Área de Tecnologia em Gestão Financeira da região, ambos avaliados com conceito 3.
Os Centros Universitários contaram com três cursos participantes na região Nordeste, um com conceito 2 e dois com conceito 4.
As Faculdades foram representadas por 13 cursos na região Nordeste.
O conceito modal foi 3, com cinco cursos.
Dos cursos restantes, um recebeu conceito 1, três receberam conceito 4 e outros três receberam conceito 5, enquanto apenas um curso ficou sem conceito ( SC ).
Na região Sudeste, as Universidades concentraram 31 dos 97 cursos da região.
Entre os cursos em Universidades na região, o conceito modal foi 3 com 14 cursos, e nenhum curso ficou sem conceito ( SC ).
Os demais cursos receberam os conceitos 2 ( oito cursos ), 4 ( seis cursos ) e 5 ( três cursos ).
Os Centros Universitários participaram com 15 cursos na região Sudeste, dos quais nove obtiveram conceito modal, 3, e nenhum ficou sem conceito ( SC ).
Os demais receberam os conceitos 2 ( cinco cursos ) e 4 ( um curso ).
As Faculdades foram representadas por 51 cursos na região Sudeste, que se distribuíram nos conceitos 1 ( três cursos ), 2 ( 12 cursos ), 3 ( 26 cursos, conceito modal ), 4 ( sete cursos ) e 5 ( dois cursos ), enquanto apenas um curso ficou sem conceito ( SC ).
Dos 37 cursos da região Sul, dez eram de Universidades, para os quais o conceito modal foi 4, com cinco cursos.
Nesse tipo de organização, nenhum dos cursos ficou sem conceito ( SC ) e os demais receberam os conceitos 2 ( um curso ), 3 ( três cursos ) e 5 ( também um curso ).
Os Centros Universitários da região Sul participaram com seis cursos, que se distribuíram nos conceitos 2 ( um curso ), 3 ( dois cursos ), 4 ( outros dois cursos ) e 5 ( também um curso ).
Foram 21 os cursos vinculados a Faculdades na região Sul, e apenas um destes ficou sem conceito ( SC ).
Cinco desses cursos receberam conceito 2, seis cursos receberam conceito 3, oito cursos receberam conceito 4, o conceito modal, e o outro curso, conceito 5.
Na região Centro-Oeste, três dos nove cursos eram de Universidades, dos quais um ficou sem conceito ( SC ), e os outros dois receberam os conceitos 3 e 5, um cada.
Os Centros Universitários da região Centro-Oeste contaram com dois cursos, dos quais um recebeu conceito 3 e o outro, conceito 4.
Nenhum curso ficou sem conceito ( SC ).
Dos quatro cursos em Faculdades na região Centro-Oeste, nenhum ficou sem conceito ( SC ), dois receberam conceito 2, um conceito 3 e o outro, conceito 4.

% matched lemmas: alocar
\end{textsample}
